\documentclass[11pt]{article}
\usepackage[a4paper,margin=25mm]{geometry}
\usepackage{amsmath,amssymb,amsthm,booktabs,array}
\usepackage[T1]{fontenc}
\usepackage[colorlinks=true,urlcolor=blue,citecolor=blue,linkcolor=blue]{hyperref}
\newtheorem{theorem}{Theorem}
\newtheorem{lemma}{Lemma}
\newtheorem{corollary}{Corollary}
\newcommand{\NN}{\mathbb N_0}
\newcommand{\ZZ}{\mathbb Z}
\newcommand{\FF}{\mathbb F_2}
\title{OEIS A323206, COMMENTS (revision 52):\\
  \large Proof of Luschny's parity conjecture and binary consequences}
\author{Haobo Ma\\\small ChronoAI Pte Ltd, Singapore
  \and Wenlin Zhang\\\small National University of Singapore\\
  \small ChronoAI Pte Ltd, Singapore}
\date{November 8, 2026}
\begin{document}
\maketitle
\begin{abstract}
We prove Peter Luschny's conjecture in OEIS A323206, COMMENTS,
revision 52 (first posed in revision 17, February 21, 2019):
for all nonnegative $n,k$, the terminating hypergeometric value $A(n,k)$ is
odd precisely when $n$ is even or $k+2$ is a positive power of two.
Its coefficients are integer ballot numbers, and their sum is $C_{k+1}$.
Thus evaluation modulo two reduces to the constant coefficient or a
Catalan number. A binary product describes the complete coefficient
support; its consequences include the linked rows, columns, three
adjacent diagonals, antidiagonal sums and generalized ballot triangle.
\end{abstract}

\section{The conjecture and its finite interpretation}
Luschny's entry \cite{posing}, contributed February 21, 2019, labels its
parity statement a ``Conjecture.'' Its first occurrence is revision 17,
COMMENTS, February 21, 2019; revision 52 retains it unchanged.
We prove that statement for every
$n,k\in\NN$, including the zero boundaries. In the NAME and COMMENTS of
revision 52 (March 15, 2026), the array and conjecture are
\begin{align}
 A(n,k)&={}_2F_1(-k,k+1;-k-1;n),\label{array}\\
 A(n,k)\text{ odd}&\iff n\text{ even}\ \text{or}\
 \bigl(n\text{ odd and }k+2=2^e\text{ for some }e>0\bigr).\label{conj}
\end{align}
Here the hypergeometric notation denotes the terminating polynomial
\begin{equation}
 P_k(t)=\sum_{j=0}^{k}
   \frac{(-k)_j(k+1)_j}{(-k-1)_j\,j!}\,t^j,
 \qquad (a)_j=\prod_{r=0}^{j-1}(a+r),\qquad A(n,k)=P_k(n).
 \label{finite}
\end{equation}
An empty product and $t^0$, including $0^0$, equal one. Every denominator
in this finite sum is nonzero: if $r<j\le k$, then $-k-1+r\ne0$.
No infinite-series interpretation at a lower-parameter pole is needed.

The ballot-number expressions in the FORMULA section of A323206 and the
Catalan triangle A009766 \cite{ballot} supply the integer structure behind
the conjecture. Catalan parity itself is classical; A000108 records both
the Catalan formula and its binary valuation \cite{catalan}.
Earlier proofs of Catalan parity include E\u{g}ecio\u{g}lu (1983)
\cite{egecioglu}; Deutsch and Sagan (2004; published 2006),
Theorem 2.1 \cite{deutsch}, give the full binary valuation.
These results concern the Catalan ingredient; the two-parameter
classification follows by the polynomial reduction below. Abate and
Whitt (2011), equations (4), (9)--(10) \cite{abate}, already give these
generalized Catalan polynomials in the notation
$C_{k+1}(t)=P_k(t)$, with the same ballot coefficients.
We give the
needed identities and their proofs below, so the parity argument is
self-contained.

\begin{theorem}\label{main}
For every $k\in\NN$, $P_k$ has integer coefficients. For every integer $t$,
\[
 P_k(t)\equiv
 \begin{cases}
 1 &\pmod2,\quad t\text{ even},\\
 C_{k+1}&\pmod2,\quad t\text{ odd},
 \end{cases}
 \qquad C_m=\frac1{m+1}\binom{2m}{m}.
\]
Consequently \eqref{conj} holds for all $n,k\in\NN$; the same criterion
holds for $P_k(t)$ at negative integers.
\end{theorem}

\section{Ballot coefficients and Catalan parity}
We take $\binom ab=0$ when $b<0$ or $b>a$, with $a\in\NN$.

\begin{lemma}\label{coeff}
For $0\le j\le k$, the coefficient of $t^j$ in $P_k(t)$ is
\begin{equation}
 b_{k,j}=\binom{k+j}{j}-\binom{k+j}{j-1}.
 \label{ballot}
\end{equation}
In particular $b_{k,0}=1$, and $\sum_{j=0}^k b_{k,j}=C_{k+1}$.
\end{lemma}
\begin{proof}
Cancellation of the nonzero factors in \eqref{finite} gives
\[
 \frac{(-k)_j}{(-k-1)_j}=\frac{k+1-j}{k+1},\qquad
 \frac{(k+1)_j}{j!}=\binom{k+j}{j}.
\]
For $j\ge1$, the ratio of $\binom{k+j}{j-1}$ to
$\binom{k+j}{j}$ is $j/(k+1)$, proving \eqref{ballot}; the case $j=0$
is immediate. Also, for $j\ge1$,
\[
 b_{k,j}=\frac{k-j+1}{j}\binom{k+j}{k+1}.
\]
Replacing $j$ by $k-j$ in this identity and in \eqref{ballot} gives
the other two finite binomial-sum formulas in A323206.
The two hockey-stick identities, obtained by telescoping
Pascal's rule, give
\[
 \sum_{j=0}^k\binom{k+j}{j}=\binom{2k+1}{k},\qquad
 \sum_{j=1}^k\binom{k+j}{j-1}=\binom{2k+1}{k-1}.
\]
Their difference is
\[
 \binom{2k+1}{k}-\binom{2k+1}{k-1}
 =\frac{2}{k+2}\binom{2k+1}{k}
 =\frac{1}{k+2}\binom{2k+2}{k+1}=C_{k+1}.
\]
The same identities include $k=0$ under the stated convention.
\end{proof}

\begin{lemma}\label{catparity}
For every $m\in\NN$, $C_m$ is odd if and only if $m+1$ is a power of two.
\end{lemma}
\begin{proof}
Write $s_2(a)$ for the number of ones in the binary expansion of $a$,
and $v_2(a)$ for the exponent of two in a positive integer $a$.
Counting the multiples of $2^r$ in $a!$ gives
\[
 v_2(a!)=\sum_{r\ge1}\left\lfloor\frac{a}{2^r}\right\rfloor
 =a-s_2(a).
\]
The second equality follows by summing the contribution $2^i-1$ of
each nonzero binary digit $2^i$. Since $s_2(2m)=s_2(m)$, the Catalan
formula yields
\[
 v_2(C_m)=s_2(m)-v_2(m+1)=s_2(m+1)-1.
\]
For the last equality, adding one to $m$ removes its $v_2(m+1)$ trailing
ones and inserts one new one. The valuation is zero exactly when $m+1$
has a single nonzero binary digit. This also covers $m=0$.
\end{proof}

\begin{proof}[Proof of Theorem~\ref{main}]
Lemma~\ref{coeff} proves integrality. In $\FF$, an even integer is zero
and an odd integer is one. Thus $P_k(t)$ reduces to $b_{k,0}=1$ in the
first case and to $\sum_j b_{k,j}=C_{k+1}$ in the second.
Lemma~\ref{catparity}, with $m=k+1$, gives $k+2=2^e$; since $k+2\ge2$,
the exponent satisfies $e>0$. These implications are equivalences,
so they prove both directions of the conjecture.
\end{proof}

\begin{corollary}[Coefficientwise binary support]\label{binary}
For $0\le j\le k$, $b_{k,j}$ is odd precisely when $j$ and $k+1$ have
no common nonzero binary digit. Equivalently,
\[
 P_k(t)\equiv\sum_{\substack{0\le j\le k\\j\mathbin{\&}(k+1)=0}}t^j
 \quad\text{in }\FF[t],
\]
where $\mathbin{\&}$ denotes binary bitwise AND.
\end{corollary}
\begin{proof}
Subtraction equals addition in $\FF$, so Pascal's rule gives
$b_{k,j}\equiv\binom{k+j+1}{j}\pmod2$.
If $N=\sum_i\epsilon_i2^i$ is a binary expansion, then
$(1+x)^N=\prod_{\epsilon_i=1}(1+x^{2^i})$ in $\FF[x]$.
Its $x^j$ coefficient is one exactly when the binary digits of $j$
are a subset of those of $N$. For $N=j+(k+1)$ this is equivalent to
addition without a carry: disjoint nonzero digits give this property;
if a common nonzero digit exists, the lowest such digit produces a zero
in the sum where $j$ has a one. This proves the criterion.
\end{proof}

\begin{corollary}[Binary factorization]\label{product}
Put $m=k+1$ and $h=\lfloor\log_2 m\rfloor$. Write $m_i$ for the $i$th
binary digit of $m$, and let $s_2(m)$ count its nonzero digits. Then
\[
 P_k(t)=\prod_{\substack{0\le i<h\\m_i=0}}(1+t^{2^i})
 \quad\text{in }\FF[t].
\]
Exactly $2^{h+1-s_2(m)}$ coefficients of $P_k$ are odd, and the degree
of its reduction modulo two is $2^{h+1}-m-1$.
\end{corollary}
\begin{proof}
Corollary~\ref{binary} allows exactly the subsets of the zero digits of
$m$ below its highest digit. Every such subset gives an index less than
$2^h\le m$, hence at most $k$; a permissible index cannot have a higher
digit. Expanding the product therefore gives precisely that support,
with no repeated exponents. There are $h+1-s_2(m)$ free digits, and their
sum is $(2^h-1)-(m-2^h)$, proving the count and degree, including $m=1$.
\end{proof}

\begin{corollary}\label{generating}
Let $C(x)=\sum_{r\ge0}C_rx^r$ and
$G(t,x)=\sum_{k\ge0}P_k(t)x^k$. In $\ZZ[t][[x]]$,
\begin{equation}
 G(t,x)=\frac{C(tx)}{1-xC(tx)},\qquad
 1+xG(t,x)=\frac1{1-xC(tx)}.\label{gf}
\end{equation}
If $U=xG(t,x)$, then
\[
 x=\frac{U(1-(t-1)U)}{(1+U)^2}.
\]
Thus $U$ is the compositional inverse of
$y(1-(t-1)y)/(1+y)^2$. Also $1+xG$ is the formal continued fraction
with first numerator $x$ and every subsequent numerator $tx$,
all subtracted from unit denominators.
\end{corollary}
\begin{proof}
Extend $b_{k,j}$ by zero outside $0\le j\le k$. Pascal's rule gives
$b_{k,j}=b_{k,j-1}+b_{k-1,j}$ for $k\ge1$, $1\le j\le k$,
and $b_{k,k}=C_k$. Hence
$(1-t)P_k=P_{k-1}-C_kt^{k+1}$. Summation over $k\ge1$ gives
$(1-t-x)G=1-tC(tx)$.
The identity $C(tx)=1+txC(tx)^2$ shows that the right side of
\eqref{gf} satisfies this equation. Multiplication by $1-t-x$ is
injective in the integral domain $\ZZ[t][[x]]$, proving \eqref{gf}
without division by $1-t$. Eliminating $C(tx)$ from
$U=xC(tx)/(1-xC(tx))$ gives the inverse equation. Its linear
coefficient is one, so the inverse is unique. Finally the tail of the
continued fraction is $C(tx)=1/(1-txC(tx))$; substitution gives
$1/(1-xC(tx))$. Finite truncations stabilize coefficientwise.
\end{proof}

\section{Consequences for the linked sequences}
Theorem~\ref{main} settles the parity of every fixed row and column:
even-$n$ rows are entirely odd, odd-$n$ rows have odd entries exactly at
$k=2^e-2$ ($e\ge1$), and a column is entirely odd precisely at those
indices $k$. In particular $A(0,k)=1$ and $A(1,k)=C_{k+1}$ by
Lemma~\ref{coeff}. The three diagonal sequences and the antidiagonal
sums linked from A323206 have the following explicit consequences
\cite{diagonals,sums}.

\begin{corollary}\label{siblings}
For $r\in\NN$, put $S_r=\sum_{k=0}^r A(r-k,k)$. Then
\begin{align*}
 A(r,r)\text{ odd}&\iff r\text{ even},\\
 A(r,r+1)\text{ odd}&\iff r\text{ even or }r=2^e-3\ (e\ge2),\\
 A(r+1,r)\text{ odd}&\iff r\text{ odd or }r=2^e-2\ (e\ge1),\\
 S_{2m}&\equiv m+1\pmod2,\\
 S_{2m+1}&\equiv m+1+\lfloor\log_2(2m+2)\rfloor\pmod2\quad(m\in\NN).
\end{align*}
\end{corollary}
\begin{proof}
Substitute the two indices into Theorem~\ref{main}. On the main diagonal,
odd $r$ makes $r+2$ odd and at least three, hence not a power of two.
The other two substitutions give the displayed conditions, including
their lower bounds on $e$.
For $S_r$, an even first index contributes one; there are
$\lfloor r/2\rfloor+1$ such indices. If $r=2m$, the odd first indices
have odd second index $k$, so none has $k+2$ a power of two.
If $r=2m+1$, the additional odd contributions are exactly
$k=2^e-2\le r-1$, or $1\le e\le\lfloor\log_2(r+1)\rfloor$.
Counting these contributions proves both congruences.
\end{proof}

\begin{corollary}[Rows, columns and the generalized Catalan triangle]
\label{rows}
For the linked sequences \cite{rows}, use the following index conventions.
Let $R_m(0)=1$ and $R_m(s)=P_{s-1}(m)$ for $s\ge1$. Then
\[
 R_m(s)\text{ odd}\iff m\text{ even or }s+1=2^e\ (e\ge1),
 \qquad(s\ge1).
\]
These are the generalized Catalan rows A064062, A064063 and
A064087--A064091 for $m=2,3,4,5,6,7,8$, respectively, with their
initial term retained. The six columns A001844, A064096,
A064302--A064305 are $P_k(n)$ for $k=2,3,4,5,6,7$:
the first and fifth are entirely odd; the other four are odd exactly
at even $n$.

In the notation $T(N,m)$ of A064094, $0\le m\le N$,
\[
 T(N,N)=1,\quad T(N,m)=P_{N-m-1}(m)\ (N>m),
\]
so for $N>m$ it is odd exactly when $m$ is even or
$N-m+1$ is a power of two.
\end{corollary}
\begin{proof}
The row entries have generating function $1/(1-xC(mx))$;
Corollary~\ref{generating} gives the stated shift. Their parity is
Theorem~\ref{main} with $k=s-1$. Substituting the six fixed $k$ values
proves the column statements. The column-$m$ generating function in
A064094 is $x^m/(1-xC(mx))$; extracting coefficients gives its displayed
identification, including $m=0$, and then the parity criterion.
\end{proof}

\begin{corollary}[Generalized ballot numbers]\label{generalballot}
Use A238762's ballot notation $B(p,q)$, $0\le p\le q$. For
$0\le j\le k$,
\[
 B(2j,2k)=b_{k,j},\qquad
 B(2j+1,2k+1)=(k+j+1)b_{k,j}.
\]
Mixed-parity arguments give zero. Consequently the even-even value is
odd exactly when $j\mathbin{\&}(k+1)=0$; the odd-odd value is odd exactly
when this condition holds and $k+j$ is even.
\end{corollary}
\begin{proof}
Substitution of the two parity choices into A238762's explicit
binomial formula gives the two displayed identities; its parity
factor is zero for mixed arguments. Apply Corollary~\ref{binary} and
reduce the factor $k+j+1$ modulo two.
\end{proof}

\begin{center}
\small
\begin{tabular}{@{}p{0.39\linewidth}p{0.55\linewidth}@{}}
\toprule
Source item & Coverage\\
\midrule
A323206, COMMENTS, revisions 17 and 52 & Proved: the quoted
conjecture, Theorem~\ref{main}.\\
A323206, the three finite binomial sums & Proved by Lemma~\ref{coeff},
  including the substitution $j\mapsto k-j$.\\
A323206, formal g.f., pseudo reversion and continued fraction &
  Proved: Corollary~\ref{generating}, including its coefficient shift.\\
A323206, EXAMPLE and CROSSREFS rows/columns; A064094 &
  Not posed; parity proved in Corollary~\ref{rows}. $n=0,1$ are exact.\\
A323209, A323208, A323217, finite sums & Proved by the corresponding
  substitutions in Lemma~\ref{coeff}; parity not posed, proved in
  Corollary~\ref{siblings}.\\
A323207, finite sums & Proved by summing Lemma~\ref{coeff};
  antidiagonal parity not posed, proved in Corollary~\ref{siblings}.\\
A009766; A238762, explicit ballot formulas & Classical formulas;
  parity not posed, proved in Corollaries~\ref{binary}, \ref{product}
  and \ref{generalballot}.\\
A000108, Catalan parity and valuation & Prior proofs cited:
  E\u{g}ecio\u{g}lu (1983), Deutsch--Sagan (2004/2006); Lemma~\ref{catparity}.\\
A323206, integral, asymptotic and Maple recurrence;
  A323208/209/217, integral and asymptotic formulas & Not covered.\\
Other formulas/questions in the row, column and ballot entries &
  Not covered beyond the explicit identifications and parity above.\\
\bottomrule
\end{tabular}
\end{center}
The linked diagonal and sum entries contain formulas, not additional
parity conjectures. The corollaries above state consequences rather than
attribute new conjectures to their authors.

\section{Mathematical verification coverage}
The pinned Lean source
\href{https://github.com/the-omega-institute/trureturing/blob/b0f7c01b01511ea00428c1b6edad5dd2b372eaf8/D5/S1/Recurrence/Invariants/LuschnyA323206Parity.lean}
{\texttt{LuschnyA323206Parity.lean}}
contains the declaration
\begin{quote}\small\ttfamily
D5.S1.Recurrence.Invariants.\\
LuschnyA323206Parity.result
\end{quote}
Its statement is \texttt{claim}: for all natural $n,k$ there is an integer
equal, after coercion to the rationals, to the finite sum \eqref{finite},
and its oddness is equivalent to \eqref{conj}.
The imported declaration
\begin{quote}\small\ttfamily
CatalanCompositionSquareParity.binary\_catalan
\end{quote}
in
\href{https://github.com/the-omega-institute/trureturing/blob/b0f7c01b01511ea00428c1b6edad5dd2b372eaf8/D5/S1/Recurrence/Invariants/CatalanCompositionSquareParity.lean}
{\texttt{CatalanCompositionSquareParity.lean}}
states that the coefficient of its shifted Catalan series at index $a$
is one modulo two exactly when $a=2^e$ for some natural $e$; this is the
Catalan support used in that proof.

The negative-integer extension and all corollaries have the written
proofs given here. The two named declarations cover the nonnegative
classification and its shifted Catalan support, respectively.
The accompanying \texttt{verify.py} performs exact finite checks of
5,151 coefficient identities and binary supports ($0\le j\le k\le100$),
101 Catalan sums, 72 printed source examples, 4,141 integer evaluations
($-20\le t\le20$, $0\le k\le100$), and 404 diagonal or antidiagonal
parity identities ($0\le r\le100$). Additional exact checks cover the
binary products, counts and degrees for $0\le k\le101$,
5,151 even-even and 5,151 odd-odd ballot values, and 10,201 mixed-parity
zeros. Formal generating and inverse identities are checked through
degree 20 at ten integer arguments; 5,151 triangular entries of A064094
are checked by its independent finite sum. These finite ranges complement the
universal written proofs.

\clearpage
\begin{thebibliography}{9}
\bibitem{posing} Peter Luschny, OEIS A323206, NAME, COMMENTS, FORMULA,
EXAMPLE and CROSSREFS; contributed February 21, 2019; revision 52,
March 15, 2026. \url{https://oeis.org/A323206}.
The conjecture first appears in revision 17, COMMENTS, February 21, 2019:
\url{https://oeis.org/history/view?seq=A323206&v=17}.
\bibitem{ballot} The OEIS Foundation Inc., A009766 (Catalan triangle),
FORMULA, and A238762 (generalized ballot numbers), FORMULA.
\url{https://oeis.org/A009766}; \url{https://oeis.org/A238762}.
\bibitem{catalan} The OEIS Foundation Inc., A000108 (Catalan numbers),
NAME and Franklin T. Adams-Watters's comment of February 8, 2006,
on the binary valuation. \url{https://oeis.org/A000108}.
\bibitem{egecioglu} \"Omer E\u{g}ecio\u{g}lu,
The parity of the Catalan numbers via lattice paths,
\emph{Fibonacci Quarterly} \textbf{21} (1983), 65--66.
\url{https://www.mathstat.dal.ca/FQ/Scanned/21-1/egecioglu.pdf}.
\bibitem{deutsch} Emeric Deutsch and Bruce E. Sagan,
Congruences for Catalan and Motzkin numbers and related sequences,
\emph{Journal of Number Theory} \textbf{117} (2006), 191--215,
Theorem 2.1; preprint July 19, 2004.
\url{https://arxiv.org/abs/math/0407326}.
\bibitem{abate} Joseph Abate and Ward Whitt,
Brownian motion and the generalized Catalan numbers,
\emph{Journal of Integer Sequences} \textbf{14} (2011), Article 11.2.6,
equations (3)--(4), (9)--(10).
\url{https://cs.uwaterloo.ca/journals/JIS/VOL14/Whitt/whitt6.html}.
\bibitem{diagonals} Peter Luschny, OEIS A323209, A323208 and A323217,
NAME and FORMULA; contributed February 25, 2019.
\url{https://oeis.org/A323209}; \url{https://oeis.org/A323208};
\url{https://oeis.org/A323217}.
\bibitem{rows} The OEIS Foundation Inc., generalized Catalan sequences
A064062, A064063, A064087--A064091; columns A001844, A064096,
A064302--A064305; triangle A064094, FORMULA and EXAMPLE.
\url{https://oeis.org/A064094}.
\bibitem{sums} Peter Luschny, OEIS A323207, NAME and FORMULA;
contributed February 25, 2019. \url{https://oeis.org/A323207}.
\end{thebibliography}
\medskip\noindent\small
\copyright\ 2026 Haobo Ma and Wenlin Zhang.
Licensed under Creative Commons Attribution 4.0 International (CC BY 4.0).
\end{document}
